\documentclass[11pt,a4paper]{article}
\usepackage{fontspec,amsmath,amssymb,amsthm,microtype}
\setmainfont{lmroman10-regular.otf}[BoldFont=lmroman10-bold.otf,ItalicFont=lmroman10-italic.otf,BoldItalicFont=lmroman10-bolditalic.otf]
\usepackage[margin=25mm]{geometry}
\usepackage[unicode,hidelinks]{hyperref}
\hypersetup{pdfauthor={LangevinLean Project}}
\setlength{\emergencystretch}{3em}
\setlength{\parskip}{0.4em}
\setlength{\parindent}{0pt}
\allowdisplaybreaks[2]
\raggedbottom
\title{\texorpdfstring{The scalar quadratic recurrence and its invariant floor}{The scalar quadratic recurrence and its invariant floor}}
\author{LangevinLean Project}
\date{}
\begin{document}\maketitle
Draft mathematical exposition; these notes may contain mistakes.


This module isolates arithmetic used after the actual LMC one-step estimate is proved. Its recurrence premise is an intermediate numerical interface, not a replacement for the stochastic theorem.



The invariant-floor calculation retains the quadratic cross term:


\[\begin{gathered}
e_{k+1}^2\le q^2e_k^2+2qCe_k+B^2,\qquad
C\le(1-q)F,\\
2qCF+B^2\le(1-q^2)F^2
\quad\Longrightarrow\quad e_n\le q^ne_0+F.
\end{gathered}\]


With the coefficients of the Laplacian theorem (Theorem 2, source label th:3), the certified floor is


\[e_n\le(1-mh)^ne_0+\frac{(2MA+\delta)hS}{m},
\qquad A=\sqrt{M_{\mathrm{av}}},\quad S=\sqrt p.\]
\section{\texorpdfstring{Keep the quadratic cross term}{Keep the quadratic cross term}}


Induct on the iteration number with the comparison \(q^k\) \(e_{0}+F\). The inequalities \(C\le (1-q)F\) and \(2qCF+B^{2}\le (1-q^{2})F^{2}\) make this comparison invariant. Nonnegativity allows comparison of squares without replacing the local error by a triangle bound.


\section{\texorpdfstring{Split the floor into two contributions}{Split the floor into two contributions}}


Choose a drift floor \(a\) and a variance floor \(b\). From \(C\le (1-q)a\) and \(B^{2}\le (1-q^{2})b^{2}\), expand the square of \(a+b\) and verify the invariant-floor hypotheses. This retains the order \(h\) accumulation of the quadratic defect.


\section{\texorpdfstring{Recover the manuscript coefficient}{Recover the manuscript coefficient}}


Set \(q=1-mh\), \(C=h^{2}(MA+\delta )S/2\) and use the manuscript-style local-defect coefficient. The conditions \(m\le M\) and \(hM\le1\) control all square roots. Explicit floors add to \((2MA+\delta)hS/m\), giving the claimed iteration bound with \(A=\sqrt{M_{\mathrm{av}}}\) and \(S=\sqrt p\).


\end{document}
